\subsection{Deriving the common polar equation}

Place the pole at the focus $F=(0,0)$ and choose the directrix
\[
D:x=d,
\qquad d>0.
\]
For a point $P=(r,\theta)$,
\[
x=r\cos\theta.
\]
Its distance from the directrix is therefore
\[
d(P,D)=d-r\cos\theta
\]
for the branch lying to the left of the directrix.

The focus--directrix condition becomes
\[
r=e(d-r\cos\theta).
\]
Collecting the terms containing $r$ gives
\[
r(1+e\cos\theta)=ed.
\]
Define the focal parameter
\[
\boxed{p=ed}.
\]
Then
\[
\boxed{r(\theta)=\frac{p}{1+e\cos\theta}}.
\]

\begin{center}
\begin{tikzpicture}[scale=0.86,>=stealth]
    \draw[axis,->] (-1.0,0)--(6.6,0) node[right] {$x$};
    \draw[gray,thick] (5.0,-3.5)--(5.0,3.5) node[above] {$x=d$};
    \coordinate (F) at (0,0);
    \coordinate (P) at (2.8,2.0);
    \coordinate (H) at (5.0,2.0);
    \draw[blue!70!black,line width=1.2pt,->] (F)--(P)
        node[midway,above left] {$r$};
    \draw[orange!85!black,line width=1.1pt] (P)--(H)
        node[midway,above] {$d-r\cos\theta$};
    \draw[helper,dashed] (P)--(2.8,0)
        node[midway,right] {$r\sin\theta$};
    \draw (1.0,0) arc[start angle=0,end angle=35.54,radius=1.0];
    \node at (1.28,0.48) {$\theta$};
    \fill[geometricSubsetColor] (F) circle (2.2pt) node[below left] {$F$};
    \fill[geometricSubsetColor] (P) circle (2.2pt) node[above left] {$P$};
\end{tikzpicture}
\end{center}

\begin{corebox}[The central idea]
The formula is not an algebraic trick. The denominator records how the distance
to the directrix changes when the viewing ray turns. Changing $e$ changes the
balance between the two distances and therefore changes the type of conic.
\end{corebox}
